\firstsection{Introduction}

\maketitle

Many data-driven graphics combine multiple visual representations to compare
data, reveal cross-view relationships, and communicate complex narratives.
Such composite visualizations are common in visual analytics and
communication~\cite{vasurvey, Survey16, deng2023survey, AI4VIS}, encompassing
patterns such as concatenation, layering, faceting, and
nesting~\cite{javed2012exploring, Composite, vaid}. 
Authoring these visualizations requires managing data, encodings,
coordinate systems, layouts, and dependencies across charts. 


Existing approaches address different stages of visualization authoring, but
chart-level composition often remains indirect. Low-level libraries~\cite{bostock2011d3} provide
fine-grained control, but require authors to manually implement and maintain
alignment, shared scales, layouts, and update logic. Interactive graphical
editors similarly offer substantial freedom in constructing and manipulating
marks and visual primitives~\cite{Lyra, DataIllustrator}. While this
mark-level flexibility supports expressive designs, authors must first
assemble primitives into coherent charts and then maintain their relationships
during composition and subsequent editing. Declarative grammars and
template-based tools facilitate common charts and compositions, but require
authors to work with predefined operators, specification structures, or
layout templates~\cite{bostock2011d3, satyanarayan2016vega, Tableau}.
Departures from these structures may require substantial specification changes
or low-level customization. These approaches leave an opportunity for
authoring abstractions that operate directly on charts while retaining
interactive control over their composition and refinement.

Realizing this opportunity raises three challenges. First, composing charts as
blocks involves more than arranging them spatially: adjacent charts may share
an axis, overlapping charts may form a semantic layer, and repeated charts may
constitute a facet. These composition relationships must remain explicit and
editable so that individual blocks and their dependencies are preserved during
refinement. Second, data adaptation becomes more complex once charts are composed. For an
individual chart, incompatibilities can often be addressed by remapping data
fields and encodings. A composition, however, introduces cross-chart
constraints: its blocks may differ in data dimensions, schemas, or levels of
granularity, while shared axes, layers, and facets impose different
compatibility requirements. The authoring process must therefore help authors
identify and reconcile such mismatches without reconstructing the constituent
charts. Third, the predefined scope of chart blocks may constrain design
exploration. A block-based approach should therefore organize blocks along
meaningful design dimensions and surface alternatives for authors to explore.
Although blocks may serve as reusable templates, they should remain
configurable and replaceable rather than prescribe fixed visual outcomes.

To inform the design of chart blocks and their alternatives, we first
conducted a formative study with visualization experts. Drawing on perceptual
principles and their design experience, the experts examined what should
constitute a reusable chart block and which alternatives should be offered for
a given design. Our analysis of their discussions identified three dimensions
for organizing chart blocks---coordinate system, visual style, and data
dimensionality. These dimensions provide a structured basis for defining block
units and surfacing design alternatives during authoring.

Guided by these findings, we developed {\system}, an interactive
authoring tool for constructing composite visualizations with reusable chart
blocks. Users drag and drop chart blocks to express spatial relationships,
which the system interprets as composition operations such as concatenation,
layering, and faceting. A data-aware recommendation engine evaluates the
dimensional compatibility of the blocks and recommends alternative chart
designs and data mappings, both to resolve mismatches and to support design
exploration. These alternatives are organized along the coordinate-system,
visual-style, and data-dimensionality dimensions identified in our formative
study. Users can then rearrange, replace, and refine individual blocks while
their composition relationships remain explicit and editable. 
Through case studies and a comparative user study, 
we demonstrate that {\system} supports the construction and iterative 
refinement of complex composite visualizations and is easier to learn 
and use than Lyra~\cite{Lyra}.





In summary, our contributions are:
\begin{itemize}
    \item A formative study that identifies coordinate system, visual style,
    and data dimensionality as key dimensions for defining reusable chart
    blocks and organizing their design alternatives;

    \item An interactive authoring tool that supports
    drag-and-drop composition with explicit and editable chart relationships,
    together with a data-aware recommendation engine for evaluating
    compatibility and suggesting chart alternatives and data mappings;

    \item Empirical findings from case studies and a comparative user study
    that characterize the capabilities, benefits, and limitations of
    block-based, data-aware authoring for composite visualizations.
\end{itemize}